\section{Static response and exposure support}\label{app:static-response}

\subsection{Direct and full responses}\label{sec:model-counterfactuals}

Starting from the strategic equilibrium
$q^0=q^{\mathrm S,*}(r_0,c_0)$, we compare responses to the same change in
opportunity values.
With rivals frozen at $q_{-i}^0$, each firm solves its new capacity problem;
the collection of these individual best responses is the direct allocation
and need not itself be an equilibrium.
Allowing rivals to respond instead gives the full allocation, the new
strategic equilibrium.
Because both calculations share the initial appeal, saturation and marginal
costs at $q^0$, full minus direct isolates rival adjustment at that origin.
The same definitions apply to linear-cost changes with opportunities fixed.
Removing rival appeal altogether gives a different benchmark: its
independently equilibrated origin changes as well as its response
(Appendix~\ref{app:independent-proof}).

\subsection{Comparative statics of rival feedback}\label{sec:model-response}

To characterize the common-origin comparison locally, first change
opportunity values while holding all cost primitives fixed.
At an equilibrium marginal contribution equals marginal cost on active
coordinates, and a small disturbance must preserve that balance.
On firm $i$'s active coordinates this requires
\begin{equation}\label{eq:firm-response}
  \underbrace{J_{ii}\,\delta q_i}_{\text{own adjustment}}
  =\underbrace{\mathcal G_i\,\delta r}
  _{\substack{\text{changed opportunity}\\\text{incentives}}}
  -\underbrace{\sum_{j\ne i}J_{ij}\,\delta q_j}
  _{\substack{\text{changed incentives}\\\text{from rivals}}}.
\end{equation}
Here $J_{ii}$ measures how quickly firm $i$'s incentive to expand declines as
it adds capacity, $J_{ij}$ measures how rival $j$'s capacity changes that
incentive, and $\mathcal G_i$ translates changes in node opportunities into
changes in the marginal returns to firm $i$'s activities.
The frozen-rival response sets every rival movement $\delta q_j$ to zero; the
full response solves all firms' equations jointly.

Write $J$ for the Jacobian of marginal cost less marginal contribution,
$J_{\mathrm d}=\operatorname{blockdiag}(J_{ii})$ for its own blocks and
$J_{\mathrm o}=J-J_{\mathrm d}$ for its rival blocks.
Let $\mathcal U$ collect positive-capacity coordinates.
Strict complementarity means that every unused activity has a strictly
negative incentive to enter, so small disturbances leave this active
set unchanged.
All inverses below are restricted to $\mathcal U$;
Appendix~\ref{app:response} supplies their entries and the projected systems.
Theorem~3 of \citet{xu_zenou_zhou_conflict_2022} gives the corresponding
active-coordinate comparative statics in their pure-cost setting under their
cost assumptions and nondegeneracy.
Here the same fixed-active-set logic supports an explicit comparison between
joint adjustment and the own-firm response with rivals frozen.

\begin{proposition}[Equilibrium response and common-origin
  feedback]\label{prop:response}
  Maintain Assumption~\ref{ass:primitives} at an equilibrium with strict
  complementarity.
  For an opportunity direction $\delta r$, the full equilibrium derivative
  $\chi\delta r$ and the direct derivative $\delta q^{\mathrm d}$ are uniquely
  determined on the active coordinates by the joint and own-block response
  systems, respectively.
  Their difference is
  \begin{equation}\label{eq:feedback}
    \chi\delta r-\delta q^{\mathrm d}
    =-J_{\mathcal U\mathcal U}^{-1}(J_{\mathrm o})_{\mathcal U\mathcal U}
    \delta q^{\mathrm d}.
  \end{equation}
\end{proposition}

Feedback transmits the direct adjustment through the rival blocks, allowing
the induced adjustments of all firms to change each firm's final response.
Appendix~\ref{app:response} proves the derivative identification.

Whereas opportunity changes alter the marginal contribution from capacity,
linear-cost changes directly alter its marginal cost.
Under strict complementarity and nearby nonnegative costs, a cost reduction
raises the marginal incentive to expand in the affected activities;
\eqref{eq:cost-response} gives its response system.
If only one firm's costs change, every other firm's frozen-rival
problem is unchanged.
Their live-rival responses then arise entirely through the recipient's
expansion and the adjustments it sets off.

The sign of feedback depends on relative appeal and activity overlap.
The rival blocks contain the difference between rival-plus-outside appeal
and own appeal, $B_{i\ell}-A_{i\ell}$, so they admit both signs
(Appendix~\ref{app:response-jacobian}).
When own appeal is smaller than rival-plus-outside appeal, rival expansion
lowers the marginal incentive to expand at that node; appeal dominance
reverses this effect, and overlapping activities combine the
node-specific effects.
Whether rivalry attenuates, amplifies or redirects a response is therefore a
property of the economy's appeal and overlap, rather than a general sign restriction.

\subsection{Scale and composition}\label{sec:model-composition}

A firm can expand without changing its business mix.
Raising capacities in two activities from $(40,60)$ to $(80,120)$ doubles
scale while leaving composition at $(0.4,0.6)$.
Raising them instead to $(60,60)$ changes the shares to $(0.5,0.5)$, although
neither activity contracts: a change in normalized mix is distinct from a
physical transfer of installed capacity.
More generally, a capacity response $h$ changes total scale by
$\mathbf1^\top h$, and at positive total capacity, differentiating normalized
composition gives
\begin{equation}\label{eq:composition-derivative}
  \frac{d}{dt}\frac{q_i+th}{\mathbf1^\top(q_i+th)}\bigg|_{t=0}
  =\frac{h-p_i\mathbf1^\top h}{S_i}.
\end{equation}
The subtraction removes proportional expansion: for $h=\alpha q_i$, the
numerator is $\alpha q_i-\alpha p_i S_i=0$.
Thus a large capacity response need not imply a large change in business mix.

For a vector $h$, write $D p_{q_i}(h)$ for the right-hand side of
\eqref{eq:composition-derivative}.
Applied to full minus direct adjustment, this derivative removes precisely
the feedback proportional to initial capacity.
An activity-separable equilibrium illustrates the remaining business-mix
response.
Under strict substitution, a marginal cost reduction in one activity expands
the recipient
there and crowds rivals out of that activity.
Positive capacity elsewhere then turns those changes into business-mix movement.
The activity-local cost shock considered here differs from the relative
opportunity news used in Section~\ref{sec:simulation}.
The result establishes business-mix movement, while movement toward or away
from competitors requires additional exposure ordering
(Appendix~\ref{app:two-site-example}).

\begin{proposition}[Proportional feedback and business-mix
  response]\label{prop:composition-feedback}
  Let $S_i>0$ and let $h_i^{\mathrm f}$ and $h_i^{\mathrm d}$ be full and
  direct capacity directions at the same initial allocation.
  Then
  \begin{equation}\label{eq:proportional-feedback}
    Dp_{q_i}(h_i^{\mathrm f})=Dp_{q_i}(h_i^{\mathrm d})
    \quad\Longleftrightarrow\quad
    h_i^{\mathrm f}-h_i^{\mathrm d}=\lambda_i q_i
    \text{ for some }\lambda_i\in\mathbb R.
  \end{equation}
  For an equilibrium interpretation, maintain an interior differentiable path
  and a fixed active set.
  In an activity-separable economy with identity demand loading and positive
  diagonal cost curvature, consider a marginal linear-cost reduction $s>0$ at
  one activity of firm $j$.
  At that activity let $A_i>0$, $D=b+\sum_k A_k$, $B_i=D-A_i$, $r,b>0$, and
  assume strict substitution $A_i<B_i$ for every firm.
  The recipient expands there and every unshocked rival contracts.
  Other activities remain unchanged.
  Each firm with positive capacity at another activity therefore
  changes business mix.
\end{proposition}

\begin{proof}[Proof idea]
  By \eqref{eq:composition-derivative}, zero normalized response means
  $h=p_i\mathbf1^\top h$, which is precisely a multiple of $q_i$.
  Apply this identity to the difference of the full and direct directions.
  For the separable benchmark, the activity Jacobian is diagonal plus rank
  one; its explicit solution has a positive recipient coordinate and negative
  rival coordinates.
  Unaffected activities solve homogeneous invertible systems.
  Appendix~\ref{app:composition-feedback} gives that solution and the
  composition argument.
\end{proof}

Attenuating a nonzero composition response changes the firm's
position even if it leaves the response's direction unchanged.

\subsection{Cash flow and competitive news}\label{sec:model-cash-flow}

The activity portfolio creates an exposure before any firm adjusts.
Write $X_{i\ell}=A_{i\ell}/D_\ell$ for firm $i$'s opportunity share, with
$D_\ell=b_\ell+\sum_j A_{j\ell}$; at fixed capacities, an opportunity change
$dr$ changes cash flow by $\sum_\ell X_{i\ell}\,dr_\ell$, so shared
opportunity exposure can generate correlated cash-flow news even before
competitive adjustment (Corollary~\ref{cor:exposure}).

The cash-flow envelope measures how much of an improvement in market
opportunities becomes additional operating cash flow once competitors respond.
Under strict complementarity, let $u=\chi\delta r$ be the equilibrium
capacity response to an opportunity direction $\delta r$, holding cost
primitives fixed, with block $u_j$ for firm $j$.
Write $k_\ell^\top$ for row $\ell$ of the activity-to-demand map $K$.
The envelope then gives
\begin{equation}\label{eq:equilibrium-cash-flow}
  \begin{split}
    \frac{d}{dt}\pi_i\bigl(q^*(r+t\delta r);r+t\delta r\bigr)\Big|_{t=0}
    ={}&\underbrace{\sum_\ell X_{i\ell}\,\delta r_\ell}
    _{\mathclap{\text{value of improved opportunities}}}\\
    &-\underbrace{\sum_\ell\frac{r_\ell A_{i\ell}}{D_\ell^2}
    \sum_{j\ne i}k_\ell^\top u_j}
    _{\mathclap{\text{effect of competitors' capacity changes}}}.
  \end{split}
\end{equation}
The first term values the disturbance at the firm's existing opportunity
shares.
The second evaluates how competitors' responses alter the contribution margin
available to the firm: $\sum_{j\ne i}k_\ell^\top u_j$ is their combined change
in appeal at node $\ell$, so their expansion lowers the value of the improved
opportunity and their contraction raises it.
At the optimizing allocation the firm's own marginal payoff is zero on
coordinates where capacity moves, so its own adjustment adds no first-order
cash-flow term.

Remark~\ref{rem:envelope} derives the split.
The payoff is stylized unlevered operating cash flow, abstracting from taxes,
working capital and depreciation; the envelope therefore supplies the operating
cash-flow component of a broader valuation, to which financing and discounting
would be added.

Holding a firm's own appeal fixed, competitors' expansion at a node lowers its
share there in proportion to the resulting growth in total appeal, so the
firm's loss depends on where it earns its margin and on how strongly a change
in margin moves its cash flow.
To express this dependence, let the share of firm $i$'s contribution margin
earned at node $\ell$ be
$\psi_{i\ell}=r_\ell X_{i\ell}/\mathcal M_i$.
Its operating leverage $\Lambda_i=\mathcal M_i/\pi_i$, contribution margin per
unit of operating cash flow, scales a change in margin into a cash-flow return.
For a capacity direction $u$, the proportional growth of total appeal at node
$\ell$ is $g_\ell=\sum_j k_\ell^\top u_j/D_\ell$, which is the same
for every firm.

\begin{proposition}[Common and differential competitive
  news]\label{prop:common-news}
  Maintain Assumption~\ref{ass:primitives} at an equilibrium with strict
  complementarity, and let $u$ be a capacity direction supported on its
  positive coordinates, such as the response to an opportunity or a
  linear-cost disturbance.
  For a firm with $\pi_i>0$, the rival term of the cash-flow change along $u$,
  expressed as a cash-flow return, is
  \begin{equation}\label{eq:rival-return}
    -\Lambda_i\sum_\ell\psi_{i\ell}\Bigl(g_\ell-\frac{k_\ell^\top
    u_i}{D_\ell}\Bigr).
  \end{equation}
  Let $\vartheta_i=\Lambda_i\psi_i$ and let
  $\zeta_i=-\vartheta_i^\top g$ denote the common-node component, before the
  own-appeal correction.
  Firms with equal operating leverage and contribution mix receive identical
  common-node news.
  For any square-integrable node-growth distribution with positive total
  common-node second moment, the retained share after removing the
  cross-sectional mean is
  $\varrho=\mathbb E\sum_i(\zeta_i-\bar\zeta)^2/
  \mathbb E\sum_i\zeta_i^2$, where $\bar\zeta=N^{-1}\sum_i\zeta_i$.
  Equal contribution mixes and leverage give $\varrho=0$.
\end{proposition}

Because firms face the same node-level growth in total appeal, differences
in their contribution mixes and operating leverage determine differences in
common-node news.
A comparison of firms within a period removes the mean of this
component $\zeta_i$, leaving the dispersion of $\vartheta_i$ in the directions
in which total appeal moves.
The own-appeal correction in~\eqref{eq:rival-return} removes each firm's own
contribution to total appeal growth, completing the rival effect; its
cross-firm variation must be retained separately.
The retained share $\varrho$ is a fraction of squared common-node response across firms,
not a fraction of the aggregate cash-flow loss.
The size of the competitive effect and the cross-sectional variation it
leaves are thus distinct: competitors' expansion can lower every firm's cash
flow substantially while barely changing how firms' cash-flow returns differ.
Appendix~\ref{app:common-news} gives the second-moment expansion and proof.
